\documentclass[11pt,letterpaper]{article}
\usepackage[T1]{fontenc}
\usepackage{lmodern}
\usepackage[margin=0.82in,headheight=14pt]{geometry}
\usepackage{amsmath,amssymb,amsthm,mathtools}
\usepackage{booktabs,array,longtable}
\usepackage{microtype,xcolor,enumitem,fancyhdr}
\usepackage[colorlinks=true,linkcolor=blue!45!black,citecolor=blue!45!black,urlcolor=blue!45!black,breaklinks=true]{hyperref}
\hypersetup{pdftitle={Direct proof certificates and smaller Diophantine consistency equations},pdfauthor={AI-assisted research for the ProveIt project},pdfsubject={Explicit Diophantine representations of ZFC refutability}}
\definecolor{ink}{RGB}{25,42,63}
\definecolor{accent}{RGB}{0,105,112}
\pagestyle{fancy}
\fancyhf{}
\fancyhead[L]{\small Direct certificates and consistency equations}
\fancyhead[R]{\small Research contribution}
\fancyfoot[C]{\thepage}
\setlength{\parskip}{4pt}
\setlength{\parindent}{0pt}
\setlist{topsep=4pt,itemsep=2pt}
\setcounter{tocdepth}{2}
\newtheorem{theorem}{Theorem}[section]
\newtheorem{lemma}[theorem]{Lemma}
\newtheorem{proposition}[theorem]{Proposition}
\newtheorem{corollary}[theorem]{Corollary}
\newcommand{\N}{\mathbb N}
\newcommand{\Z}{\mathbb Z}
\newcommand{\Con}{\operatorname{Con}}
\newcommand{\ZFC}{\mathrm{ZFC}}
\newcommand{\file}[1]{\nolinkurl{#1}}
\setlength{\emergencystretch}{2em}
\newcommand{\pin}{\texttt{08200eb10fde2db408c69dffd840d6920a569b1f}}
\title{\color{ink}\LARGE Direct proof certificates and smaller\\Diophantine consistency equations\\[7pt]
\large A materialized quartic in 3,614 natural unknowns\\
and equations in 21, 20, and 19 integers\\
of degrees 16, 18, and 40}
\author{Prepared for Vladimir Reshetnikov's ProveIt project\\
\small AI-assisted research contribution with reproducible exact algebra audits}
\date{24 October 2026\\[3pt]\small Repository snapshot: \pin}
\begin{document}
\maketitle
\begin{abstract}
We give two complementary improvements to explicit Diophantine
representations of ZFC refutability. First, we materialize a fixed quartic
in 3,614 nonnegative integer unknowns, with 15,996 nonzero monomials and
maximum absolute coefficient 127,778,545. Its construction compiles the
repository's calibrated numerical proof calculus directly, without
simulating a Turing machine or instantiating Jones's universal polynomial.
A companion quartic has 6,183 unknowns and coefficient height 1,993.
The central compilation step is a shared Chinese-remainder certificate:
one global factorial makes every local range factorial invertible, allowing
the latter to be removed and heterogeneous witness bounds to be used.
The exported equations contain only ordinary polynomial monomials.

Second, integral triangular coordinate changes and two Pell-root
translations turn an existing 25-integer octic into fixed equations in
21, 20, and 19 integers of exact degrees 16, 18, and 40. At these dimensions
the best inspected incoming degrees were 20, 32, and 88.
All new coordinate maps are polynomial bijections on zero sets.
Composing the arguments with the explicitly identified inherited
calibration theorems gives, in ZFC, equivalence between absence of solutions
and $\Con(\ZFC)$. This is an ordinary mathematical proof with a stated
dependency boundary, not a new end-to-end proof-assistant certificate.
Neither consistency nor global minimality is claimed.
\end{abstract}
\newpage
\begingroup
\small
\setlength{\parskip}{0pt}
\tableofcontents
\endgroup
\newpage

\section{The precise result and the meaning of simplicity}

Fix the usual arithmetized assertion $\Con(\ZFC)$ that no finite ZFC proof
of contradiction exists. The target is a fixed integer-coefficient
polynomial $P$, with its domain specified, for which
\begin{equation}\label{eq:target}
 \ZFC\vdash
 \bigl(\neg\exists\boldsymbol{x}\in D^n\ P(\boldsymbol{x})=0\bigr)
 \ \longleftrightarrow\ \Con(\ZFC).
\end{equation}
Here $D$ is either $\N=\{0,1,\ldots\}$ or $\Z$. These are different
representation problems: a natural-number equation must not silently be
advertised as an unrestricted integer equation. All coefficients are fixed;
there is no free input encoding the theory.

There is no single established numerical meaning of the simplest
Diophantine equation. Variable count, degree, coefficient height, number
of monomials, circuit size, readability, and the complexity of the
calibration proof compete. We therefore deliver several fully specified
points, rather than conceal a large cost in an unspecified coefficient
or arithmetic macro.

\begin{table}[h]
\centering\small
\begin{tabular}{lrrrr}
\toprule
Equation and domain & Unknowns & Degree & Squared residuals & Coefficient information\\
\midrule
$P_{\mathrm{dir}}$ over $\N$ & 3,614 & 4 & 2,954 & height 127,778,545\\
$P_{\mathrm{ht}}$ over $\N$ & 6,183 & 4 & 5,523 & height 1,993\\
$Q_{21}$ over $\Z$ & 21 & 16 & 9 & six supplied numerals\\
$Q_{20}$ over $\Z$ & 20 & 18 & 8 & same six numerals\\
$Q_{19}$ over $\Z$ & 19 & 40 & 7 & same six numerals\\
\bottomrule
\end{tabular}
\caption{Completed equations. Height means the maximum absolute coefficient
of the expanded polynomial, not of a residual or a circuit.}
\end{table}

\begin{theorem}[Completed consistency representations]\label{thm:main}
The five fixed polynomials in the table satisfy \eqref{eq:target} on
their stated domains, by the ordinary arithmetic arguments below composed
with the repository calibration results specified in
Section~\ref{sec:scope}. Their variable counts and degrees are exact.
\end{theorem}

The direct construction is useful when small coefficients, degree four,
and a proof-calculus interpretation matter. The 19-integer construction
is the most compact new equation here by variable count; it retains the
large fixed numerical tuple of its inherited machine-based source.
The two routes have different calibration dependencies.

\subsection{What is new relative to the inspected repository}

All comparisons refer to the snapshot \pin, including the six archive
roots in \file{docs/incoming}. The accompanying comparison inventory
identifies each source. In particular, the incoming joint-guard package
already contains a 25-integer octic and several small-variable projections.
Reconstructing those equations is necessary calibration work, not a new
result. The new coordinate projections improve the inspected
$(n,\deg)$ points
\[
 (21,20)\mapsto(21,16),\quad
 (20,32)\mapsto(20,18),\quad
 (19,88)\mapsto(19,40).
\]
No claim is made to improve existing nine-positive or ten-integer
high-degree representations by variable count.

On the direct route, the repository already supplies a numerical Horn
calculus, commutative proof certificates, and a bounded-factor carrier.
Its large-carrier candidate reports 47,439,839 representation witnesses
before a final ordinary-polynomial flattening. That number is a different
compilation stage, so a ratio against our final 3,614 coordinates would
not be a controlled like-for-like comparison. Our contribution is the
shared compilation theorem, the combination of concrete reductions,
and a completed literal polynomial with small numerical coefficients.
We found no such completed small-coefficient direct quartic among the
inspected incoming packages. This is a repository comparison, not a claim
of worldwide priority.

\subsection{An unambiguous static definition of the direct polynomial}

The file \file{direct/equation.json} lists 2,954 residuals. A residual is
a finite list of pairs $(c,\mu)$, where $c\in\Z$ and $\mu$ is a list of
zero, one, or two variable indices in $\{0,\ldots,3613\}$. Repetition of
an index means a square. The empty list means the constant monomial.
If the residual lists are $\mathcal L_0,\ldots,\mathcal L_{2953}$, the
single equation is, exactly,
\begin{equation}\label{eq:literal-polynomial}
 P_{\mathrm{dir}}(X_0,\ldots,X_{3613})
 =
 \sum_{\ell=0}^{2953}
 \left(\sum_{(c,\mu)\in\mathcal L_\ell}
          c\prod_{j\in\mu}X_j\right)^2=0.
\end{equation}
These lists are included as data, not left as the output of an unavailable
compiler. The fully expanded coefficient list is also included in
\file{direct/expanded_polynomial.json}. The builder reproduces both.
The analogous two files in \file{direct/small_height} define
$P_{\mathrm{ht}}$. No exponentiation, binomial coefficient, factorial,
division, bounded quantifier, or variable-length product remains in these
files. Formula \eqref{eq:literal-polynomial} is simply a compact convention
for displaying a completely fixed polynomial.

\section{The exact numerical calculus encoded by the source}

The local source uses the numerical calculus $C_{222}$, independently of
any machine that enumerates it.  Its primitive codes are the following
ordinary nonnegative polynomials:
\begin{align*}
 I(a,b)&=2(a+b)(a+b+1)+4b,\\
 M(a,b)&=I(a,b)+1,& A(a,b)&=I(a,b)+2,& E(a,b)&=I(a,b)+3,\\
 N(a)&=I(a,1),& X(a,b)&=N(A(a,N(b))).
\end{align*}
They are, respectively, the repository's implication, membership,
universal quantification, equality, negation, and existential-quantification
codes.  The falsity code is $1$, and the free truth code is $0$.

The eight unconditional axiom codes are
\begin{align*}
 r_1(a,b)&=X(a,E(a,b)),\\
 r_2(a,b,c)&=I(A(a,I(b,c)),I(A(a,b),A(a,c))),\\
 r_3(a,b,c)&=I(E(a,b),I(M(c,a),M(c,b))),\\
 r_4(a,b,c)&=I(A(a,A(b,c)),A(b,A(a,c))),\\
 r_5(a,b,c)&=I(E(c,a),I(M(c,b),M(a,b))),\\
 r_6(a,b)&=I(X(a,A(a,b)),b),\\
 r_G(a,b,c)&=I(E(c,a),I(E(c,b),A(M(a,b),E(a,b)))),\\
 r_8(a,b,c,d)&=I(I(I(c,d),a),I(I(a,c),I(b,c))).
\end{align*}
The guarded seed bundle is also a fixed composition of these polynomials.
To state it without a long nested line, define
\begin{align*}
 R(b,k,h)&=A(b,I(h,N(A(b,I(h,N(M(b,k))))))),\\
 S(a,b,k)&=A(k,I(M(k,a),M(k,b))),\\
 C_0(a,b,k,h)&=I(S(a,b,k),R(b,k,h)),\\
 C_1(a,b,k,h)&=N(A(a,C_0(a,b,k,h))),\\
 C_2(a,b,k,h)&=N(A(a,I(h,C_1(a,b,k,h)))),\\
 C(a,b,k,h)&=A(b,I(M(b,k),C_2(a,b,k,h))),\\
 V(a,b,c)&=I(M(a,b),A(c,M(a,b))),\\
 K(a,b,c)&=A(c,I(C(a,b,c,c),N(R(b,c,c)))),\\
 \mathcal B(a,b,c)&=N(I(V(a,b,c),K(a,b,c))).
\end{align*}
The bundle is an axiom when $c\ne a$ and $c\ne b$.  These conditions
are represented using one natural witness $d$ by
\[
 (c-a)^2(c-b)^2-(d+1)^2=0.
\]
Existence of such $d$ is equivalent to the two disequalities: in the
forward direction the product cannot be zero, and in the reverse direction
take $d=|(c-a)(c-b)|-1$.

There are precisely twelve branches in the padded two-premise relation:
truth $(p,q,v)=(0,0,0)$; each of the eight unconditional axiom instances
with $p=q=0$; the guarded bundle with $p=q=0$; modus ponens
$q=I(p,v)$; and generalization $q=0$, $v=A(a,p)$.
The four natural schema coordinates $a,b,c,d$ are shared across this
finite disjunction.  Their values need satisfy only the chosen branch.

\paragraph{An exact arithmetic audit, beyond random replay.}
The supplementary script \nolinkurl{audit_calculus.py} independently
implements the displayed formulas and compares them with the delivered
source by sparse-polynomial arithmetic.  A composite expression is replaced
by a source wire name only after its coefficient dictionary equals that
wire's defining polynomial exactly.  In this way it recognizes all fifty-six
branch-specific wire definitions and verifies all thirty-four remaining
source residuals.  Each replacement is a finite ring identity.  This
checks the arithmetic translation compositionally, without expanding the
very large nested bundle polynomial and without treating a finite random
sample as an identity proof.

The direct comparison is against
\nolinkurl{Code254.lean} and \nolinkurl{KCode222.lean} at the pinned repository
commit.  The semantic dependencies are then
\[
 C_{222}(1)
 \ \longleftrightarrow\
 \operatorname{BProv}(\mathrm{ZF},\bot)
 \ \longleftrightarrow\
 \operatorname{BProv}(\mathrm{ZFC},\bot).
\]
The named repository theorems are
\nolinkurl{KEndpoint.c222_iff_zf_refutation} and
\nolinkurl{ConstructibleChoice.zf_refutation_iff_zfc_refutation}.
The latter uses the constructible-universe proof that ZF and ZFC are
equiconsistent.  The current construction inherits those mathematical
calibration results.  It does not claim to have rechecked their full Lean
dependency closure, and the new compiler is not itself a Lean proof.


\section{Why a finite commutative certificate is enough}
\label{sec:flow}

This section records the mathematical content of the inherited direct
flow mechanism~\cite{flow}, so that the bridge used by the compiler is
visible. Let $\mathcal H$ be a finite list of instances of
single-conclusion Horn rules with at most two premises. Axiom instances
have zero premises. Pad missing premises with the freely available truth
label $0$, and let $e_0=0$. For a rule $(p,q)\to v$, assign the flow
vector $e_v-e_p-e_q$ in the free abelian group on the nonzero labels.

\begin{lemma}[Single-conclusion flow]
A nonzero label $b$ is derivable if and only if a finite multiset of
rule instances has sum of flow vectors $e_b$.
\end{lemma}
\begin{proof}
A finite proof can be expanded to a finite derivation tree. Every internal
formula occurrence cancels between its production and its use; the root
contributes $e_b$, and truth leaves contribute zero.

Conversely let $D$ be the least set of labels containing truth and closed
under the rules, and suppose $b\notin D$. Sum the coordinates outside
$D$ of each rule vector. A rule with conclusion outside $D$ has at least
one premise outside $D$, by closure, so its contribution is at most zero.
A rule with conclusion in $D$ also contributes at most zero.
Every nonnegative integral sum of rule vectors therefore has nonpositive
total outside $D$, whereas $e_b$ has total one. This is a contradiction.
Only the finitely many labels in the proposed certificate are needed
for the finite sum.
\end{proof}

\subsection{A carrier that avoids pairwise-prime formula weights}

Take $t\ge1$, $m\ge1$, with every displayed formula label at most $m$,
and define
\[
 M=(m+1)^t,\qquad w(a)=1+aM.
\]
For $t$ rule instances with conclusions $v_i$ and premises $p_i,q_i$,
consider
\[
 A(X)=\prod_{i<t}(1+v_iX),\qquad
 C(X)=(1+X)\prod_{i<t}(1+p_iX)(1+q_iX).
\]
Suppose $A(M)=C(M)$. First $\deg C\le t$: otherwise
$C(M)\ge M^{t+1}$, while
\[
 A(M)\le(1+mM)^t<((m+1)M)^t=M^{t+1}.
\]
Each of $A$ and $C$ now has at most $t$ nonconstant linear factors, all
with labels at most $m$. Its coefficient sum is at most
$(m+1)^t=M$. Its constant coefficient is one, so every coefficient is
strictly less than $M$. Uniqueness of ordinary base-$M$ expansion implies
$A=C$ in $\Z[X]$. Unique factorization over $\mathbb Q[X]$ then gives
the flow equation with target $1$: each positive label contributes its
own distinct factor $1+aX$. The free label zero contributes the unit.
No coprimality of the numerical weights $1+aM$ is assumed.

\subsection{An accumulator with a polynomial local step}

Introduce natural states $S_0,\ldots,S_t$ satisfying
\[
 S_0=M+1,\quad S_t=1,\qquad
 S_{i+1}(1+v_iM)=S_i(1+p_iM)(1+q_iM).
\]
Every state is positive, by induction from $S_0>0$, since all weights
are positive. Multiplying the local equations and cancelling the
intermediate states gives $A(M)=C(M)$, hence a proof of the target.

Conversely expand a target derivation as a finite tree and traverse it
by replacing a frontier conclusion by its premises. Let $S_i$ be the
product of the weights on that frontier. This supplies natural accumulator
states satisfying every step, starting with $1+M$ and ending with $1$.
Choose $t$ to be the number of rule occurrences and $m$ to bound the labels,
then choose the carrier $M$ by the displayed formula.

The states are encoded by one ordinary beta sequence
\[
 \beta(x,y,i)=x\bmod(1+(i+1)y).
\]
Choose $y$ greater than every required state and divisible by $t!$.
The moduli for $0\le i\le t$ are pairwise coprime by the elementary
gcd argument used below, so finite CRT supplies $x$ with the prescribed
digits. Crucially, this beta-sequence scale $y$, the carrier $M$, and
the later shared-compiler scale $T$ are chosen separately. Identifying
these scales would introduce an unjustified size cycle.

\section{A shared CRT compiler with binomial bounds}
\label{sec:shared-crt}

All variables in this section range over $\mathbb N=\{0,1,2,\ldots\}$.
Let $E(k,z,w)$ and $F(k,z,w)$ be fixed polynomials with nonnegative
integer coefficients, where $k$ is one index, $z$ is a finite tuple of
external parameters, and $w=(w_1,\ldots,w_d)$ is a finite witness tuple.
The nonnegative coefficient hypothesis is used only for monotonicity.
For a proposed vector of witness bounds $u=(u_1,\ldots,u_d)$, choose a
natural $V$ with $u_j\le V$ for every $j$, and put
\[
 H=E(N,z,u)+F(N,z,u)+N+V+1,\qquad T=H!,\qquad
 B=\prod_{k=0}^{N-1}(1+(k+1)T).
\]
The empty product is $1$.  In particular $T\ge1$, $B\ge1$, and
$H\ge u_j+1$ for every coordinate.

\begin{theorem}[Shared binomial certificate]
\label{thm:shared-binomial}
For fixed $N,z,u,V$ as above, the following assertions are equivalent:
\begin{enumerate}
\item For every $k<N$ there is $w\in\mathbb N^d$ with
      $w_j\le u_j$ for all $j$ and $E(k,z,w)=F(k,z,w)$.
\item There are $C,A_1,\ldots,A_d\in\mathbb N$ such that
\begin{align}
 B&\mid\binom{A_j}{u_j+1} &&(1\le j\le d),\label{eq:choose-bounds}\\
 B&\mid1+(C+1)T,\label{eq:index-divisibility}\\
 E(C,z,A)&\equiv F(C,z,A)\pmod B.\label{eq:evaluation-congruence}
\end{align}
\end{enumerate}
\end{theorem}

\begin{proof}
Write $D_k=1+(k+1)T$.  We have $\gcd(D_k,T)=1$.
The numbers $D_0,\ldots,D_{N-1}$ are pairwise coprime: if $i<j<N$
and $g=\gcd(D_i,D_j)$, then $g\mid(j-i)T$ and $\gcd(g,T)=1$,
so $g\mid j-i$.  But $j-i\le N\le H$, hence $j-i\mid H!=T$,
and $g=1$.  Consequently $\gcd(B,T)=1$ as well.

Suppose first that the arithmetic certificate exists.  Fix $k<N$.
Choose a prime $p\mid D_k$; this is possible since $D_k>1$.
Then $p\mid B$ and $p\nmid T$.  A prime at most $H$ divides $H!$,
so $p>H$.  For every $j$, the integer identity
\[
 (u_j+1)!\binom{A_j}{u_j+1}
       =\prod_{r=0}^{u_j}(A_j-r)
\]
and \eqref{eq:choose-bounds} show that $p$ divides the product on the
right.  This identity also holds when $A_j\le u_j$, because both sides
are zero.  Primality therefore supplies $w_j\in\{0,\ldots,u_j\}$
with $A_j\equiv w_j\pmod p$.  Comparing
\eqref{eq:index-divisibility} with $D_k\equiv0\pmod p$, and cancelling
$T$ modulo $p$, gives $C\equiv k\pmod p$.  Polynomial evaluation and
\eqref{eq:evaluation-congruence} imply
\[
 E(k,z,w)\equiv F(k,z,w)\pmod p.
\]
Monotonicity gives
\[
 0\le E(k,z,w)\le E(N,z,u)<p,
 \qquad 0\le F(k,z,w)\le F(N,z,u)<p.
\]
These congruent natural numbers are equal, proving the bounded row.

Conversely, choose one bounded witness $w^{(k)}$ for each $k<N$.
By the Chinese remainder theorem and the pairwise coprimality just proved,
choose natural $C,A_j$ satisfying
\[
 C\equiv k\pmod{D_k},\qquad
 A_j\equiv w_j^{(k)}\pmod{D_k}
 \quad(k<N,\ 1\le j\le d).
\]
The latter congruence makes the falling-factorial product vanish modulo
$D_k$, since its factor indexed by $w_j^{(k)}$ vanishes.  Thus
\[
 B\mid (u_j+1)!\binom{A_j}{u_j+1}.
\]
Now $(u_j+1)!$ divides $H!=T$ and $\gcd(B,T)=1$; we may cancel this
factor to obtain \eqref{eq:choose-bounds}.  The index congruences give
\eqref{eq:index-divisibility}; the valid rows and polynomial congruence
give \eqref{eq:evaluation-congruence}.  Pairwise coprimality combines
the congruences modulo the $D_k$ into congruences modulo $B$.
When $N=0$, there are no row obligations and $B=1$, so both directions
hold without any separate side condition.
\end{proof}

\paragraph{Why the factorial cancellation matters.}
The repository's supplied CRT certificate uses
$B\mid(u_j+1)!\binom{A_j}{u_j+1}$ separately for every coordinate.
Its exact source does not require $H\ge u_j+1$, so simply dropping these
factorials from that source would be unjustified.  The explicit added bound
$V+1$ repairs this issue.  It also permits heterogeneous bounds at no
additional factorial cost.  A representation of each binomial coefficient
is still necessary; the theorem does not eliminate that obligation.

\subsection{A finite arithmetic expression for the modulus product}

The product defining $B$ is explanatory notation and is eliminated before
the final polynomial is issued.  Set
\[
 Q=(1+(N+1)T)^N+2,\qquad L=TQ-1.
\]
Then the exact natural-number identity is
\begin{equation}
 \label{eq:product-expression}
 B=\left(T^N N!\binom{Q+N}{N}\right)\bmod L.
\end{equation}
Indeed, the numerator equals
$\prod_{i=1}^N(TQ+iT)$ and reduces modulo $L$ to
$\prod_{i=1}^N(1+iT)=B$.  Further,
\[
 B\le(1+NT)^N<(1+(N+1)T)^N+1\le TQ-1=L,
\]
with the strict inequality also valid for $N=0$, when the left side is
$1$ and the middle upper bound is $2$.  Thus $B$ is the remainder.
The denominator is positive because $T\ge1$ and $Q\ge3$.
This is the same explicit product mechanism used by the repository's
\nolinkurl{BoundedForallRepresentation.productGraph}; the elementary
calculation above records the needed identity independently.

Powers, binomial coefficients, factorials, and the final remainder are each
replaced by their fixed finite polynomial graph.  In particular, no
variable-length product, variable exponent, or bounded universal quantifier
remains in the emitted ordinary polynomial.

\subsection{Keeping the majorants as arithmetic circuits}

A signed arithmetic circuit $P$ admits natural polynomial circuits $E,F$
such that $P=E-F$.  The recursive pairs are $(X,0)$ for a variable,
$(c,0)$ or $(0,-c)$ for an integer constant according to its sign,
\[
 (E_1,F_1)+(E_2,F_2)=(E_1+E_2,F_1+F_2),
\]
\[
 (E_1,F_1)-(E_2,F_2)=(E_1+F_2,F_1+E_2),
\]
and
\[
 (E_1,F_1)(E_2,F_2)=
 (E_1E_2+F_1F_2,\ E_1F_2+F_1E_2).
\]
Every operation preserves nonnegative coefficients.  Repeated subexpressions
are shared as nodes in a finite directed acyclic graph.  Expanding this DAG
into a tree is mathematically harmless but can make the representation far
larger.  The compiler uses node values as natural witnesses and emits their
quadratic defining equations only once per use-site evaluation.

\section{Eliminating local Horn wires before bounded compilation}

The pinned \nolinkurl{affine_shared_step.json} source contains twelve Horn
branches.  Its existential coordinates include thirteen shared names
$q_0,\ldots,q_{12}$.  In each branch the actual wire definitions are
triangular polynomials with nonnegative coefficients, and the first rows of
that branch assert exactly $q_i-f_i=0$.  Consequently, on natural inputs,
the definitions determine a unique natural value for every wire used by that
branch.  Substitution removes all thirteen existential wire coordinates.

There are fifty-six branch-specific definition occurrences and thirty-four
remaining branch residuals.  These are verified from the pinned JSON by exact
sparse-polynomial comparison.  Branch definitions use separate scopes.  This
does not assume one set of wire values makes all branches vanish: existential
quantification distributes over the finite disjunction, and precisely one
vanishing branch is sufficient.  The product of the twelve branch sums of
squares implements that disjunction.

The affine step originally has twenty-three natural witness coordinates.
Wire elimination leaves ten:
\[
 p,q,v,a,b,c,d,\delta_p,\delta_q,\delta_v.
\]
Here $(p,q)\to v$ is the rule, $a,b,c,d$ are its four schema witnesses,
and the three deltas implement $p,q,v\le m$.
Using Theorem~\ref{thm:shared-binomial} with heterogeneous bounds removes
the latter three witnesses too: bound $p,q,v$ directly by $m$ in the CRT
certificate.  Bound all remaining per-row witnesses by $V=m+U$, where
$U$ is one additional outer natural variable.  Such $U$ always exists for
any finite valid certificate.

\subsection{Thirteen per-row coordinates}

Retain the five outer parameters $t,x,y,m,M_0$ and put $M=M_0+1$.
The carrier condition is $M=(m+1)^t$.  At row $i<t$, set
$D_i=1+(i+1)y$ as an arithmetic expression, without introducing a
separately bounded row coordinate.  Besides the seven Horn coordinates
$p,q,v,a,b,c,d$, use six row coordinates $S,S',u,r,u',r'$.
The four beta equations are
\begin{align*}
 x-S-D_i u&=0,& S+r+1-D_i&=0,\\
 x-S'-(D_i+y)u'&=0,& S'+r'+1-(D_i+y)&=0.
\end{align*}
Since $D_i,D_i+y\ge1$, these say exactly that $S$ and $S'$ are
the consecutive beta digits.  The remaining accumulator equation is
\[
 S'(1+vM)=S(1+pM)(1+qM).
\]
Let $P$ be the product of the twelve substituted Horn branch sums of squares,
plus the square of this accumulator residual and the squares of the four
beta residuals.  It is a fixed nonnegative polynomial in the index,
the external parameters and these thirteen row coordinates.  Split its
arithmetic circuit as $P=E-F$ as above.

For the first three encoded coordinates use the bound $m$, and for the
other ten use $V=m+U$.  Form
\[
 H=E(t,z,u)+F(t,z,u)+t+V+1,
 \qquad u=(m,m,m,V,\ldots,V),
\]
where $z=(t,x,y,m,M_0)$ denotes the external parameters (the polynomial
need not actually use $t$ except through the index bound).
The shared certificate then expresses all rows with thirteen binomial
graphs and no per-coordinate factorial graphs.

The endpoints need four natural witnesses $u_0,r_0,u_t,r_t$ and the rows
\begin{align*}
 x-(M_0+2)-(y+1)u_0&=0,\\
 M_0+r_0+2-y&=0,\\
 x-1-(1+(t+1)y)u_t&=0,\\
 r_t+2-(1+(t+1)y)&=0.
\end{align*}
They assert $\beta(0)=M+1$ and $\beta(t)=1$ exactly.

\paragraph{Logical dependency boundary.}
The changes in this section preserve the numerical C222 Horn rule predicate
of package 68.  They do not supply a new proof that this predicate is
equivalent to ZFC refutability.  That calibration is inherited from the
repository's pinned C222-to-ZF and constructible-choice arguments, including
the sealed target code $1$ and free truth code $0$.  The equivalence of
complete accumulator certificates with finite Horn proofs is the pinned
single-conclusion flow theorem, composed with the bounded-factor carrier
theorem.  Given those dependencies, all subsequent transformations here are
explicit arithmetic equivalences provable in ZFC; none assumes
$\operatorname{Con}(\mathrm{ZFC})$.

The pertinent sources are pinned to the repository snapshot on the
title page:
\begin{itemize}
\item The DiophantineConsistency Lean modules for bounded-forall
      certificates and their representation, Pell equations,
      exponentiation, and factorials.
\item Package 68's \nolinkurl{COMMUTATIVE_CERTIFICATES.md} and
      \nolinkurl{affine_shared_step.json}.
\item The candidate \nolinkurl{direct-flow-large-carrier/README.md}.
\end{itemize}
The provenance inventory supplies their complete repository paths and
the immediate source copies.

\section{Specializing every power graph to positive arguments}

The completed source can choose $t\ge1$ and $m\ge1$ without losing a
certificate.  A zero-length accumulator path contradicts its distinct
endpoints, and if all labels were zero then every accumulator update would
leave the state unchanged.  Hence the existence predicate is preserved by
the parameterizations $t=t_0+1$, $m=m_0+1$.

Likewise each encoded CRT coordinate may be required positive.  The Chinese
remainder theorem specifies residue classes only; adding a positive multiple
of $B$ preserves all required congruences.  Thus one may parameterize
$A_j=\alpha_j+1$ for natural $\alpha_j$.  This affects encoded coordinates,
not the decoded row witnesses: the latter may still be zero, including the
free truth code.  The range certificate remains valid because its proof
already establishes divisibility of the binomial value for every CRT
representative of the prescribed small residues.

Under these choices every power graph used in the construction has positive
base and positive exponent.  The carrier and modulus powers have exponent
$t>0$; the factorial graphs are called at $t$ and $H$, both positive;
their radius and auxiliary binomial arguments are positive.  In each
coordinate-range binomial graph the upper argument $A_j$ is positive and
the lower argument $u_j+1$ is positive.  Therefore its three power calls
$2^{A_j}$, $(2^{A_j}+2)^{A_j}$, and $(2^{A_j}+1)^{u_j+1}$ are positive
cases.  The zero-base and zero-exponent branches of a general power graph
are unnecessary in this specialized construction.  Their removal is
justified by this argument inventory; it is not a correct optimization of
the unrestricted exponentiation relation.

For the positive-power certificate, its final Pell norm gives
\[
 A^2=1+W^3(W+2)Z^2,
 \qquad W\ge n\ge1,\quad A\ge2.
\]
Consequently $Z\ge1$ and $A>W\ge n$.  Hence the natural subtraction
$A\mathbin{\dot{-}}n$ in the original power theorem equals ordinary
integer subtraction.  The congruence can be written with nonnegative
polynomial sides as
\[
 X+Yn+Tq_+=YA+m+Tq_-.
\]
This is the only subtraction-elimination point that needs a separate
inequality argument in the power specialization.

\section{Eliminating every arithmetic macro}
\label{sec:primitives}

This section specifies the fixed finite graph certificates used by the
direct materialization. All coordinates here are nonnegative integers.
A displayed product or square is subsequently expressed by addition and
multiplication gates; the final gate rows are quadratic.

\subsection{A positive-power graph with 27 auxiliary coordinates}

For positive inputs $n,k$ and proposed output $m$, write
\[
 \mathcal N(a,x,y)=x^2+y^2-1-a^2y^2.
\]
The positive-power graph has the following 27 additional coordinates:
\[
\begin{gathered}
 W,A,T,Z,X,Y,U,V,S,R,B,\ a_0,y_0,b_0,q,d,\\
 b_+,b_-,s_+,s_-,r_+,r_-,x_+,x_-,g,w_n,w_k.
\end{gathered}
\]
Impose the finite equations
\begin{align}
 A&=a_0+2,&Y&=k+y_0,&B&=b_0+2,\label{eq:powbounds}\\
 \mathcal N(A,X,Y)&=0,&
 \mathcal N(A,U,V)&=0,&
 \mathcal N(B,S,R)&=0,\label{eq:pownorms}\\
 B&=1+4Yq,& V&=Y^2(d+1),\label{eq:powdiv}\\
 B+Ub_+&=A+Ub_-,&
 S+Us_+&=X+Us_-,&
 R+4Yr_+&=k+4Yr_-,\label{eq:powcong}\\
 X+Yn+Tx_+&=YA+m+Tx_-,\label{eq:powfinalcong}\\
 2An&=T+n^2+1,& T&=m+g+1,\label{eq:powmodulus}\\
 W&=n+w_n,&W&=k+w_k,&
 \mathcal N(W+1,A,WZ)&=0.\label{eq:powlast}
\end{align}
Existence of these coordinates is equivalent to $m=n^k$.
This is the positive specialization of the Pell--Matiyasevich
power criterion~\cite{pell,carneiro}, with every strict inequality
represented by a natural slack, every divisibility represented by a
quotient, and every possibly signed congruence quotient represented as
a difference of two naturals. The first three norms and congruences
characterize the $k$th Pell solution at parameter $A$. The last norm,
the bounds on $W$, the modulus definition, and the strict output bound
are precisely the remaining power criterion.

The condition $Y\ge k\ge1$ excludes the exceptional zero-index Pell
solution. The modulus $4Y$ is positive. Also \eqref{eq:powlast} gives
\[
 A^2=1+W^3(W+2)Z^2.
\]
Since $A\ge2$ and $W\ge n\ge1$, one has $Z\ge1$ and $A>W\ge n$.
Thus the integer rearrangement in \eqref{eq:powfinalcong} is equivalent
to the original natural-subtraction congruence involving $A-n$.
These side conditions are part of the proof, not optional simplifications.

The optional unrestricted builder retains the zero-exponent and zero-base
branches. The delivered constructions do not need those branches, because
the preceding section and an independently checked call-site induction
establish positive inputs to all 55 power uses.

\subsection{Binomial coefficients by a base digit}

For $n,k\in\N$, put
\[
 q=2^n+1,\qquad L=q^k,\qquad P=(q+1)^n.
\]
The binomial theorem expands $P$ in base $q$, with the coefficient
$\binom nk$ in digit position $k$, because all coefficients are at most
$2^n<q$. In particular, positions above $n$ have digit zero.
For proposed output $c$, add just four natural coordinates $h,r,r_s,c_s$
besides the three power graphs, and impose
\begin{equation}\label{eq:choosegraph}
 P=L(qh+c)+r,\qquad L=r+r_s+1,\qquad q=c+c_s+1.
\end{equation}
The last two equations say $0\le r<L$ and $0\le c<q$. Uniqueness of
quotient, remainder, and base digit proves that \eqref{eq:choosegraph}
holds exactly when $c=\binom nk$.

\subsection{Factorials as a ceiling quotient}

For $n\in\N$, define
\[
 \rho=(n+2)^{n+2},\qquad c=\binom{\rho+n}{n},\qquad P=(\rho+1)^n.
\]
Then
\begin{equation}\label{eq:factorialformula}
 n!=\left\lceil\frac Pc\right\rceil.
\end{equation}
For $n=0,1$ this is immediate. For $n\ge2$,
\[
 \frac Pc=n!\prod_{i=1}^n\frac{\rho+1}{\rho+i}
 =n!\prod_{i=1}^n\left(1-\frac{i-1}{\rho+i}\right).
\]
Using $1-\prod(1-a_i)\le\sum a_i$ for $0\le a_i\le1$, we obtain
\[
 0\le n!-\frac Pc
 <\frac{n!\,n(n-1)}{2(\rho+1)}<1.
\]
The final strict bound follows, for example, from
$n!\,n(n-1)/2<n^{n+2}<(n+2)^{n+2}=\rho$.
This proves \eqref{eq:factorialformula}, including the case where the
quotient is already integral.

For proposed output $f$ and two natural coordinates $r,s$, impose
\[
 P+c=cf+r+1,\qquad c=r+s+1.
\]
Since $c>0$, these say $c(f-1)<P\le cf$, which is equivalent to
$f=\lceil P/c\rceil$. The required powers and binomial graph are those
already specified.

\subsection{Ordinary quadratic rows and one quartic}

Every arithmetic-circuit addition or multiplication is represented by a
new natural coordinate $z$ and the row
\[
 z-x-y=0\qquad\hbox{or}\qquad z-xy=0.
\]
Numerals and shared nodes are treated once at each distinct evaluation.
Signed equality is kept as an equality between two nonnegative
circuits. Each graph obligation is then an ordinary residual of degree
at most two. A finite conjunction of rows $R_\ell=0$ is equivalent over
$\N$ to the one equation $\sum_\ell R_\ell^2=0$.

Before affine elimination, the resulting fixed system has 6,183
coordinates and 5,523 rows. Every residual coefficient has absolute
value at most four. Its expanded squared sum has 17,430 monomials and
coefficient height 1,993. These are measured exact coefficient
dictionaries, rather than an asymptotic bound.

\subsection{Removing nonnegative affine gates}

Suppose a retained gate has the defining row $z-L=0$, where $L$ is an
affine polynomial with nonnegative integer coefficients in earlier
coordinates. Substituting $L$ everywhere and deleting this defining
row and coordinate is a bijection on natural zero sets: the inverse
restores $z=L$. Substitution of an affine polynomial preserves the
quadratic degree bound.

We apply this operation in topological order when the expanded affine
replacement has at most four terms. The support threshold controls
coefficient and expression growth; it is unrelated to correctness.
Exactly 2,569 coordinates and rows disappear. This yields
\[
 6183-2569=3614,\qquad 5523-2569=2954.
\]
The resulting quartic has 15,996 monomials and coefficient height
127,778,545, which is less than $2^{27}$. The increase in coefficient
height is the explicit cost of removing these affine wires.

\begin{table}[h]
\centering
\begin{tabular}{lrr}
\toprule
Measured quantity & Retained wires & Affine elimination\\
\midrule
Outer coordinates & 5 & 5\\
Other witness coordinates & 1,665 & 1,665\\
Gate coordinates & 4,513 & 1,944\\
Power / binomial / factorial graphs & 55 / 16 / 2 & 55 / 16 / 2\\
Quadratic rows & 5,523 & 2,954\\
Expanded monomials & 17,430 & 15,996\\
Maximum residual coefficient & 4 & 4,556\\
Maximum quartic coefficient & 1,993 & 127,778,545\\
\bottomrule
\end{tabular}
\caption{Exactly measured direct materializations of the same certificate.}
\end{table}

\section{Coordinated changes of variables: 21, 20, and 19 integers}
\label{sec:coordinated}

The incoming \emph{Joint Sign Guards} package already supplies a
25-integer octic and equations in 21, 20, and 19 integers of degrees
20, 32, and 88. Thus reproducing the octic is a calibration step, not
a new frontier claim. The following coordinated changes of variables
give new equations in the same respective dimensions of exact degrees
\textbf{16, 18, and 40}. They use the same complete source-34 numerical
tuple and the same two determinant scales. Every transformation below
is an integer-polynomial bijection on zero sets after an integral
triangular change of coordinates.

\subsection{The fixed coefficients and the common abbreviations}

The fixed coefficients are
\[
 \beta=2^{1064921}-1,\qquad K_0,M_C,M_F,\qquad
 A=2^{1113},\quad B=2^{666104}.
\]
Here $K_0,M_C,M_F$ are the fully supplied source-34 integers;
$M_F$ is already shifted. The six numerals are bound by the packaged
coefficient manifest and reconstructed from the literal machine/table
recipe. They occupy 615883 minimal unsigned bytes. They are fixed
coefficients, never existential parameters. In particular $\beta>0$.

For the 21-variable presentation take
\[
\begin{gathered}
 J_3,C,z,f,h,r,v,w_3,s,k,y,\gamma,o,j,Y,c,I,t,p_1,p_2,p_3
 \quad\in\mathbb Z.
\end{gathered}
\]
The letter $s$ is the shifted first Pell root (\texttt{tau\_gap} in the
literal packets); it has no connection with the scale variable $v$.
Define polynomial abbreviations
\begin{gather*}
 J_0=3J_3+1,\qquad W=4(3w_3+1),\qquad q=\beta J_0+1,\qquad X=Wq,\\
 F=(K_0+W)C-z(q-1),\qquad E=XY,\qquad a=(X+1)Y,\\
 \eta=c-Yk,\qquad\zeta=k-\eta,\qquad V=fo-c,\\
 \Delta=a^2+4a+3,\qquad D=X+ac+\gamma(4a+3),\\
 R=(q-F)(q^3-q)-C(q^2-1)+(M_C+qM_F)J_0,\\
 G=AFBC-(q-AF-BC)^2,\qquad Z=\eta\zeta,\\
 T=1+p_1^2+p_2^2+p_3(p_3+1).
\end{gather*}
The alignment is typographical only: all names on the left are
polynomial abbreviations and add no unknowns.

The nine residuals of the 21-variable equation are
\begin{align*}
 R_1&=Y-(2v+1)q^3,\\
 R_2&=s^2+(2s-k)EYk-1,\\
 R_3&=D^2-\Delta c^2-1,\\
 R_4&=k-1-hE-R,\\
 R_5&=V-jc+k-1-hE,\\
 R_6&=I-(2r+1)c^2,\\
 R_7&=I^2-\Delta(f^2-1),\\
 R_8&=y^2-1-I^2(y^2-V^2),\\
 R_9&=GZ-(t^2-G-Z+1)^2-T.
\end{align*}
Set $Q_{21}=\sum_{i=1}^9R_i^2$. The displayed degrees are,
in this order,
\[
 4,6,8,5,4,3,8,6,8.
\]
In the highest-degree part of $Z=(c-Yk)(k-c+Yk)$, the monomial
$-Y^2k^2$ occurs with coefficient $-1$. The coefficient of $Y^4k^4$
in $R_9$ is consequently $-1$, and no other residual contributes
$Y^8k^8$ to its square. Thus
\[
 [Y^8k^8]Q_{21}=1,
\]
and the exact degree is sixteen, including after all fixed ports are
specialized to their actual values.

\subsection{Why these coordinates are effective}

The literal incoming 25-integer packet uses
$k_0=1-V+jc+hE$ in three arithmetic rows, while retaining
$\eta+\zeta$ in its quotient and positivity rows. To obtain the convenient
uniform presentation used here, put
$\varepsilon=V-jc+\eta+\zeta-1-hE$. Add $(E-a)\varepsilon$ to its main-ordinate row,
add $-E(c-\eta)\varepsilon$ to its first-norm row, and add $\varepsilon$ to its packed-rank
row. Keep $\varepsilon$ itself. These are reversible elementary row operations,
so the integer zero set is unchanged. The exact checker compares all
thirteen resulting rows with the independently reconstructed packet.

In this presentation of the source 25-integer octic use the integral triangular changes
\[
 C=q-F-\alpha,\qquad Y=a-E,\qquad k=\eta+\zeta,
 \qquad s=\tau-E(c-\eta).
\]
Their inverses are
\[
 \alpha=q-F-C,\quad a=E+Y,\quad\zeta=k-\eta,
 \quad\tau=s+E(c-\eta).
\]
The source rows then determine four coordinates monically:
\[
 F=(K_0+W)C-z(q-1),\qquad E=XY,\qquad
 \eta=c-Yk,\qquad V=fo-c.
\]
Delete the four defining rows and substitute these four integer
polynomials everywhere. This reduces 25 coordinates and 13 residuals
to 21 coordinates and 9 residuals. The first norm is controlled by
the exact identity, with $u=E(c-\eta)=EYk$,
\[
 (s+u)^2-(k+u)^2+(k+u)k-1
   =s^2+(2s-k)u-1.
\]
The cancellation prevents a needless squaring of the degree-five
expression $EYk$.

Retaining both $Y$ and $k$ is essential to the bound. It keeps
$\eta,\zeta$ quadratic, even after $E$ and $F$ have been removed.
The determinant $G$ and product $Z$ are both quartic, so the combined
positivity residual remains of degree eight. Expanding the quotient
coordinates too early was responsible for the larger degree of the
inspected incoming projections.

For complete rigor, map an old zero forward by the four triangular
definitions and then forget $F,E,\eta,V$. The four deleted source
equations imply precisely the restoration polynomials above. Conversely,
restore these coordinates at any new zero and invert the triangular
changes; every source residual vanishes, including the four deleted
defining equations. The two maps are inverse on zero sets. There is
no division, denominator condition, or sign assumption in this argument.
Therefore the inherited equivalence with finite ZFC refutability transfers
to $Q_{21}$.

\subsection{Twenty integers and degree eighteen}

In $Q_{21}$ make the integral change $y=fo+b$, where $b$ is unrestricted,
and eliminate $I=(2r+1)c^2$ using $R_6$. All other coordinates remain.
The key cancellation is
\[
 (fo+b)^2-(fo-c)^2=(c+b)(2fo+b-c).
\]
The resulting auxiliary residual is
\[
 (fo+b)^2-1-(2r+1)^2c^4(c+b)(2fo+b-c),
\]
of degree nine. The seven other residuals have degrees at most eight.
The final equation $Q_{20}=0$ is their sum of squares, has twenty integer
coordinates, and has exact degree eighteen because
\[
 [r^4c^{10}f^2o^2]Q_{20}=64.
\]
Only the auxiliary row contributes this monomial. The variable change and
monic elimination are reversible over the integers. Hence $Q_{20}$ has a
zero exactly when $Q_{21}$ does.

\subsection{Nineteen integers and degree forty}

Finally substitute
\[
 Y=(2v+1)\bigl(\beta(3J_3+1)+1\bigr)^3
\]
in $Q_{20}$ and remove $R_1$. This leaves the nineteen coordinates
\[
 J_3,C,z,f,h,r,v,w_3,s,k,b,\gamma,o,j,c,t,p_1,p_2,p_3.
\]
Use exactly the common abbreviations above, the additional definitions
$I=(2r+1)c^2$ and $y=fo+b$, and the seven residuals
$R_2,R_3,R_4,R_5,R_7,R_8,R_9$. These are completely specified ordinary
polynomials, with respective degrees
\[
 12,14,7,7,14,9,20.
\]
Their squared sum is $Q_{19}$. Its exact degree is forty, since
\[
 [v^8J_3^{24}k^8]Q_{19}=2^8\,3^{24}\,\beta^{24}>0.
\]
This coefficient is the image of the $Y^8k^8$ certificate. Every other
residual has degree at most fourteen before squaring, so it cannot alter
the degree-forty coefficient. The monic restoration of $Y$ again gives a
bijection of integer zero sets.

\subsection{Proof boundary and measured comparison}

Each new system therefore has a zero exactly when the existing
25-integer octic has a zero. Combining with its named native/compiler,
signed-Pell, determinant and K215/refutation calibration gives, by ordinary
mathematics formalizable in ZFC,
\[
 \mathrm{ZFC}\ \vdash\quad
 \neg\exists\boldsymbol{x}\in\mathbb Z^n\,
       Q_n(\boldsymbol{x})=0
 \quad\longleftrightarrow\quad\operatorname{Con}(\mathrm{ZFC}),
 \qquad n\in\{19,20,21\}.
\]
No absence of zeros or consistency assertion is proved here. No new
end-to-end proof-assistant certificate or internally encoded derivation
in the repository's object-language ZFC calculus is claimed.

The comparison includes all six incoming archives visible at the pinned
repository snapshot, not just numbered proposals. The changes in
$(n,\deg)$ are
\[
 (21,20)\longrightarrow(21,16),\qquad
 (20,32)\longrightarrow(20,18),\qquad
 (19,88)\longrightarrow(19,40).
\]
These are degree reductions of 20\%,
43.75\%, and approximately 54.55\% at unchanged dimensions. Other
measures may increase; fixed coefficient payload size is unchanged.
Existing ten-integer and nine-positive high-degree equations remain
different tradeoffs. No worldwide priority or global minimum is asserted.

Exact finite checks verify every defining-row identity, both integral
root translations, all compact-display residuals, all variable counts and
liveness, the stated residual degree bounds, and the three nonzero
coefficient certificates. They also reconstruct the complete numerical
tuple. These computations audit the literal algebra; the zero-set maps
above supply the universal proof.


\section{Logical scope, semantic dependencies, and exact audits}
\label{sec:scope}

\subsection{Composing the direct proof}

The direct direction from a hypothetical zero to a refutation is:
\begin{enumerate}
\item Restore the eliminated nonnegative affine gates. A sum of squares
      vanishes precisely when every residual does.
\item Read each arithmetic gate and each positive-power, binomial,
      factorial, and remainder certificate. Their graph theorems give
      their intended exact natural values.
\item Apply the shared binomial certificate to obtain all $t$ valid
      bounded local rows. The beta equations and endpoint equations
      produce one complete accumulator path.
\item Apply the carrier and single-conclusion flow arguments to obtain
      a finite $C_{222}$ derivation of the sealed target $1$.
\item Apply the pinned C222-to-ZF and constructible-choice calibrations
      to obtain a finite ZFC refutation.
\end{enumerate}
Every implication reverses existentially: a refutation gives a finite
calculus proof; its expanded tree supplies a path; finite bounds and CRT
supply the shared witnesses; the arithmetic graph existence theorems
supply all auxiliary coordinates; and gates are evaluated. Positive
representatives and positive carrier parameters are available as proved
above. Consequently
\[
 \exists\boldsymbol{x}\in\N^{3614}\,
 P_{\mathrm{dir}}(\boldsymbol{x})=0
 \quad\longleftrightarrow\quad
 \exists p\ \operatorname{Proof}_{\ZFC}(p,\bot).
\]
The companion equation is linked by explicit gate projection and
restoration. Taking negations proves the required equivalences.

\subsection{What is inherited, and what has been established here}

The mathematical dependencies are explicit:
\begin{enumerate}
\item The numerical code constructors and axiom/rule predicates from
      \file{Code254.lean} and \file{KCode222.lean}.
\item The repository statements
      \file{KEndpoint.c222_iff_zf_refutation} and
      \file{ConstructibleChoice.zf_refutation_iff_zfc_refutation}.
      They connect this concrete calculus to the intended theory.
\item The Pell-index characterization and power graph, represented in
      Mathlib's \file{PellMatiyasevic} development and in the pinned
      repository Pell and power modules~\cite{pell,carneiro}.
\item For the integer family, the older native/compiler, signed-Pell,
      determinant-guard, exact-coefficient, and K215/refutation calibration
      of the 25-integer parent. Immediate sources are included in
      \file{coordinate_family/upstream}.
\end{enumerate}
The new work proves the shared compiler theorem, supplies elementary
graph arguments, explains the direct flow/carrier bridge, proves all new
coordinate maps, and materializes and independently audits the actual
finite polynomial data. It does not rebuild the transitive proof-assistant
dependency closure of the inherited sources. It does not supply a new
object-language ZFC proof term or claim that the Python programs establish
logical soundness by themselves.

The proofs use finite induction, finite CRT, integer arithmetic and the
cited number-theoretic theorems. These are provable in ZFC and apply to
its internal natural numbers. There is no step requiring arithmetic
soundness of ZFC, a model of ZFC, a standardness assumption about all
proof codes, or $\Con(\ZFC)$ as a hypothesis. The calibration theorems
are mathematical dependencies to be checked on their merits, not a
consistency assumption. This distinction is essential when expressing
\eqref{eq:target} as a theorem of ZFC.

\subsection{Independent finite algebra checks}

The direct checker uses a separate exact implementation to expand all
2,954 residual squares and compare the resulting coefficient dictionary
with the delivered 15,996-term polynomial. It checks canonical indices,
degree four, all variable occurrences, and the coefficient height. It
also recovers all 2,569 eliminated affine definitions from the original
rows and final variable names, without using the exporter's substitution
map, and verifies the full restored row multiset.

The arithmetic controls include complete Pell certificates for six
positive power examples, zero-exponent and zero-base controls for the
optional total graph, changed-output rejection, 63 binomial cases
including $k>n$, 20 factorial controls, 30 modulus-product controls,
and a nontrivial shared CRT fixture. A static induction proves positivity
at all 55 power call sites. These finite fixtures are regression checks;
the universal validity comes from the mathematical graph theorems.

A separate compositional sparse-polynomial audit recognizes all 56
source wire definitions and all 34 branch residuals from the displayed
numerical calculus. Branch fixtures and modular evaluations provide
additional implementation controls. A separate symbolic power-graph
audit substitutes all 61 definitional gates and matches all 17 remaining
constraints to independently written versions of the cited Pell and
power criteria, with exactly 27 additional natural witnesses.

For the integer family, another independently implemented parser and
sparse polynomial engine imports none of the author's polynomial engine.
It checks the exact row bridge from the literal incoming parent, all four
triangular inverse identities, every retained and deleted row after each
projection, variable liveness, degree bounds, and the three nonzero
coefficient certificates. The full fixed tuple is also reconstructed
from its finite recipe. Both audit families pass in ordinary Python and
with assertions disabled using \file{python3 -O}.

These checks do not search for a zero of the complete equations and do
not claim that the equations have no zeros. Their purpose is to close
the gap between the displayed algebra and the supplied literal artifacts.

\section{Artifacts and offline reproduction}

The archive is designed as a self-contained contribution directory:
\begin{description}[style=nextline]
\item[\file{report/}]
This source, its included sections, the compiled PDF, and a shell build
script. The article contains the ordinary proofs and states their scope.
\item[\file{direct/}]
The two literal direct equations, fully expanded coefficient lists,
deterministic generator, exact arithmetic and source audits, receipts,
and immediate calibration/source provenance.
\item[\file{coordinate_family/}]
The three new integer equations, their compact parent, exact six-numeral
payloads and reconstruction recipe, proofs, inherited source material,
comparison inventory, and independent checkers.
\item[\file{verify.py}]
A root verification entry point. It checks packaged hashes and runs the
independent algebra audits; its help lists optional reproduction checks.
\item[\file{MANIFEST.sha256}]
Hashes of the packaged deliverables, excluding the manifest itself.
\end{description}

The numerical and polynomial audits use Python's standard library.
Typical commands from the extracted directory are
\begin{verbatim}
python3 verify.py
python3 verify.py --help
cd report
sh build.sh
\end{verbatim}
The document build uses a normal pdfLaTeX installation with the packages
listed in the preamble. The JSON files are deliberately plain and
inspectable; the mathematical definitions do not require trusting a
particular executable serialization format. The hash manifest protects
the identity of the files audited, not the truth of their mathematics.

\section{What would constitute the next substantial advance?}

The strongest conceptual gain here is reuse of one large arithmetic
certificate across many local obligations. The factorial-free range
step exposes several concrete remaining questions.

\paragraph{Share more than the global factorial.}
The thirteen coordinate bounds still require thirteen binomial graphs.
The direct materialization contains 55 power certificates, so improvement
of a shared multiple-binomial or shared-power graph would affect the
largest repeated component. A correct construction must preserve each
coordinate's decoded residue; simply combining vanishing products does
not establish all coordinate bounds.

\paragraph{Replace the factorial scale by a cheaper sufficiently divisible scale.}
The shared theorem needs the relevant small differences and range
factorials to divide $T$, and it needs every prime factor of each
$1+(i+1)T$ to exceed the evaluation bounds. The factorial is a transparent
way to guarantee all three. A smaller graph for a scale with these
properties could help, but a numerical reduction in $T$ alone need not
reduce the number of polynomial unknowns.

\paragraph{Avoid duplicate evaluation circuits.}
The compiler evaluates a nonnegative majorant and a residue instance of
the same large numerical rule circuit. A direct congruence circuit might
share additional structure, but it must still provide an honest bound
for prime decoding. Integer cancellation may make an expression smaller
while destroying the nonnegative-coefficient monotonicity used in the
proof.

\paragraph{Optimize the calibration rather than the simulated machine.}
A smaller directly calibrated proof calculus could reduce both circuit
evaluations and the branch product. Such a proposal must include a
complete equivalence with ZFC refutability. A short uncalibrated calculus,
or an inconsistent strengthening, would not solve the stated problem.

\paragraph{Continue the integer coordinate search with degree-aware translations.}
The 21-variable construction succeeds because it keeps $Y=a-E$ and
$k=\eta+\zeta$ as coordinates while eliminating other apparently equally
natural choices. The first root is shifted before expansion; the second
root is shifted before its companion coordinate is removed. A systematic
search over integral triangular coordinates could uncover further
cancellations. Exact zero-set inverses and a nonzero leading coefficient
should be required for every candidate.

\paragraph{Measure several objectives together.}
The direct quartic with height 1,993 is larger than the variant with
height 127,778,545. Both are meaningful outcomes of the same compiler.
An optimization report should retain the full tuple of domain, unknowns,
degree, height, monomials and source-circuit size. A change that is excellent
for one objective may be poor for another.

These are specific avenues opened by the present constructions. None is
being counted as a completed further improvement. Establishing an absolute
minimum is a separate problem, and the deliverable makes no such claim.

\clearpage
\begin{thebibliography}{9}
\raggedright
\bibitem{repo}
V. Reshetnikov, \emph{ProveIt}, repository snapshot
\pin.\\
\url{https://github.com/VladimirReshetnikov/ProveIt}.
Relevant directories: \file{Computability/HilbertTenthProblem},
\file{Computability/DiophantineConsistency/con-zfc-diophantine-forms},
and \file{docs/incoming}.

\bibitem{flow}
ProveIt, package 68, \emph{Commutative certificates} and the
\emph{direct-flow-large-carrier} candidate, at the same snapshot.
Immediate source text and the exact affine step are included with this
contribution.

\bibitem{pell}
The Mathlib community, \emph{PellMatiyasevic}: Pell-index
Diophantine representation and power characterization.
\url{https://leanprover-community.github.io/mathlib4_docs/Mathlib/NumberTheory/PellMatiyasevic.html}.
The relevant names include \file{Pell.matiyasevic} and
\file{Pell.eq_pow_of_pell}. The immediate repository adaptations are
\file{PellRepresentation.lean} and \file{PowerRepresentation.lean}.

\bibitem{carneiro}
M. Carneiro, \emph{A Lean formalization of Matiyasevi\v{c}'s theorem},
2018. \url{https://arxiv.org/abs/1802.01795}.

\bibitem{incoming}
ProveIt incoming proposals at the pinned snapshot:
\emph{Joint Sign Guards}; \emph{Structural Reductions, Round 11};
\emph{Budget and Centering}; \emph{Centered Pell Norms};
\emph{Radical Guard Descent}; and \file{report.zip}.
The complete archive-name/hash inventory and inspected root texts are
included in the coordinate-family comparison and provenance files.

\bibitem{sun}
Z.-W. Sun, \emph{Ten unknowns for Hilbert's tenth problem over the
integers}, 2026. \url{https://arxiv.org/abs/2609.01594}.
This provides context for small-variable consistency equations;
the constructions here are defined and proved through the dependencies
explicitly stated above.
\end{thebibliography}
\end{document}
